%% kapa: exam=1 round=36 year=2025 period=1 form=A subject=real_estate subject_ko=부동산학원론 seq=2 number=82
%% source: Q-Net 감정평가사 기출문제 (cst00309) artlSeq=5231525 file=raw/1008/5231525/2025년 제36회 감정평가사 1차시험 1교시 문제지 원본.pdf
%% conversion: kordoc 4.24.51 -> tools/gen1.py (auto); review=proofread-sonnet; flags=none
\begin{question}{82}
준부동산(準不動産)에 관한 설명으로 옳지 않은 것은?
\begin{choices}
  \item 준부동산은 특정의 부동산, 동산과 부동산의 집단을 말한다.
  \item 공장재단이란 공장에 속한 일정한 기업용 재산으로 구성되는 일단(一團)의 기업재산으로 ｢공장 및 광업재단 저당법｣에 따라 소유권과 저당권의 목적이 되는 것을 말한다.
  \item 광업재단이란 광업권과 광업권에 기하여 광물을 채굴ㆍ취득하기 위한 각종 설비 및 이에 부속하는 사업의 설비로 구성되는 일단의 기업재산으로 ｢공장 및 광업재단 저당법｣에 따라 소유권과 저당권의 목적이 되는 것을 말한다.
  \item 입목이란 토지에 부착된 수목의 집단으로서 그 소유자가 ｢입목에 관한 법률｣에 따라 소유권보존등기를 받은 것을 말하며, 토지소유권 또는 지상권 처분의 효력이 미친다.
  \item ｢수산업법｣의 규정에 따른 어업권이란 면허를 받아 어업을 경영할 수 있는 권리를 말한다.
\end{choices}
\end{question}
\answer{1,4}
\answernote{복수 정답(정답표 원문: 1, 4)}
